% Generated draft: requires mathematical and source review.
\documentclass[11pt]{article}
\usepackage[margin=25mm]{geometry}
\usepackage{fontspec}
\setmainfont{Latin Modern Roman}
\usepackage{amsmath,amssymb,mathtools,unicode-math}
\setmathfont{Latin Modern Math}
\usepackage{enumitem,needspace,xurl,hyperref}
\setlength{\parindent}{0pt}
\setlength{\parskip}{5pt}
\setlength{\emergencystretch}{3em}
\newcommand{\N}{\mathbb{N}}
\newcommand{\Z}{\mathbb{Z}}
\newcommand{\Q}{\mathbb{Q}}
\newcommand{\R}{\mathbb{R}}
\newcommand{\C}{\mathbb{C}}
\newcommand{\F}{\mathbb{F}}
\newcommand{\A}{\mathbb{A}}
\newcommand{\PP}{\mathbb{P}}
\newcommand{\E}{\mathbb{E}}
\newcommand{\eps}{\varepsilon}
\newcommand{\dd}{\,\mathrm{d}}
\begin{document}
\section*{UnsolvedMath ID: 8800094}
\textbf{Source problem code:} catalogue entry
\subsection{Definitions and mathematical statement}
The source requests explicit test elements or principal homogeneous spaces with prescribed ramification and prescribed large prime order \(\ell\), but does not identify their base field, acting group, cohomology group or local conditions.

For a commutative algebraic group \(A\) over a field \(K\), a principal homogeneous space is a torsor under \(A\), represented by a class \([T]\in H^1(K,A)\) in the appropriate topology. In this abelian-group setting, exact prime order means
\[
 \ell[T]=0,\qquad[T]\ne0.
\]
For a noncommutative acting group \(H^1\) is generally only a pointed set, so ``order'' itself would need another meaning. No particular elliptic curve, torus or finite Galois module is supplied.

Likewise ramification requires a specified local cohomological model. For example, for a finite Galois module over a global field, an unramified local class is normally tested by restriction to inertia at a place. A Selmer-group formulation also requires its module or torsion level and chosen local subgroups. These notions are contextual possibilities, not definitions recovered from the excerpt.

The intended construction must specify admissible ramification data and any global compatibility conditions, what \(\ell\) is allowed to depend on, and what explicit output is required, such as equations or cocycles. The word ``prescribed'' is not promoted to a claim that every arbitrary collection of local conditions is realizable. Nor is a polynomial-time bound inserted merely from ``explicit''.

The summary interprets the problem as explicit Selmer/class-field-theoretic construction and reports no verified general solution, but supplies no exact construction theorem or missing objects. The 2006 request and later open assessment remain source attributions pending primary-literature review.
\subsection{Short English statement}
Construct the intended torsors or test classes with exact prime order \(\ell\) and specified ramification. The base field, acting group, local conditions and compatibility requirements are missing, so no universal realization theorem is assumed.
\subsection{Sources}
\begin{itemize}
\item Local registry: \url{data/unsolvedmath/problems.json}, permanent UnsolvedMath ID \(8800094\), source code catalogue entry; supplied statement, background, and research summary, updated 2026-08-06. Attributed to Ming-Deh Huang.
\item IPAM, \emph{Workshop on Number Theory and Cryptography---Open Problems} (2006), attachment 6673, PDF page 30: \url{https://www.ipam.ucla.edu/programs/workshops/workshop-i-number-theory-and-cryptography-open-problems/?tab=schedule}. Supplied PDF/legacy-PPT visual audit accessed 2026-07-29; it reports 19 inspected surviving attachments.
\item No separate definition of the test objects, Selmer conditions or exact construction/resolution reference is supplied.
\end{itemize}
\end{document}
